\documentclass[11pt]{article}
\usepackage[a4paper,margin=1in]{geometry}
\usepackage{times}
\usepackage[T1]{fontenc}
\usepackage[utf8]{inputenc}
\usepackage{microtype}
\usepackage{booktabs}
\usepackage{setspace}
\usepackage[hidelinks]{hyperref}
\usepackage{url}
\usepackage{hanging}
\emergencystretch=3em

% Every measurement in this article comes from here, generated by
% make_numbers.py from the analysis outputs. Nothing is typed.
% Generated by make_numbers.py from the analysis outputs.
% Do not edit: every value here has exactly one source, and it is
% not this file. Regenerated before each build.

\newcommand{\nAffirmed}{10{,}021}
\newcommand{\nAffirmedAppealProc}{95.2}
\newcommand{\nAffirmedFirstProc}{92.6}
\newcommand{\nAffirmedShare}{65.1}
\newcommand{\nAppealInstruments}{29}
\newcommand{\nAppealOnly}{1{,}459}
\newcommand{\nAppealProcedural}{95.6}
\newcommand{\nAppeals}{15{,}383}
\newcommand{\nBothReasoned}{3{,}106}
\newcommand{\nBothReasonedShare}{22.3}
\newcommand{\nDecidedOnMerits}{11{,}238}
\newcommand{\nDiffAffirmed}{$-$1.2}
\newcommand{\nDiffAll}{$-$1.5}
\newcommand{\nDiffDisturbed}{$-$7.1}
\newcommand{\nDisturbedAppealProc}{89.2}
\newcommand{\nDisturbedFirstProc}{93.7}
\newcommand{\nDisturbedShare}{10.8}
\newcommand{\nFirstInstruments}{42}
\newcommand{\nFirstProcedural}{93.2}
\newcommand{\nHolmCitations}{1.000}
\newcommand{\nHolmDammam}{.029}
\newcommand{\nHolmDefence}{.778}
\newcommand{\nHolmInstruments}{1.000}
\newcommand{\nHolmJeddah}{1.000}
\newcommand{\nHolmLength}{.307}
\newcommand{\nHolmMerits}{.157}
\newcommand{\nHolmProcedural}{.006}
\newcommand{\nHolmReasons}{.699}
\newcommand{\nHolmRiyadh}{1.000}
\newcommand{\nHolmSurvivors}{three}
\newcommand{\nHolmYear}{<.001}
\newcommand{\nJudgments}{50{,}666}
\newcommand{\nModelBase}{11.8}
\newcommand{\nModelFullBase}{10.8}
\newcommand{\nModelFullN}{11{,}238}
\newcommand{\nModelN}{9{,}563}
\newcommand{\nNoReasonsShare}{74.4}
\newcommand{\nNotAdmitted}{2{,}261}
\newcommand{\nNotAdmittedShare}{14.7}
\newcommand{\nOrCitations}{0.98}
\newcommand{\nOrDammam}{1.44}
\newcommand{\nOrDefence}{1.13}
\newcommand{\nOrInstruments}{1.05}
\newcommand{\nOrJeddah}{1.09}
\newcommand{\nOrLength}{1.11}
\newcommand{\nOrMerits}{0.86}
\newcommand{\nOrProcedural}{0.74}
\newcommand{\nOrReasons}{0.85}
\newcommand{\nOrRiyadh}{1.09}
\newcommand{\nOrYear}{1.30}
\newcommand{\nOtherDisposition}{435}
\newcommand{\nOtherShare}{2.83}
\newcommand{\nOwnReasons}{3{,}931}
\newcommand{\nOwnReasonsShare}{25.6}
\newcommand{\nPCitations}{.509}
\newcommand{\nPDammam}{.003}
\newcommand{\nPDefence}{.156}
\newcommand{\nPDiffAffirmed}{.280}
\newcommand{\nPDiffAll}{.087}
\newcommand{\nPDiffDisturbed}{.033}
\newcommand{\nPInstruments}{.485}
\newcommand{\nPJeddah}{.438}
\newcommand{\nPLength}{.044}
\newcommand{\nPMerits}{.020}
\newcommand{\nPProcedural}{<.001}
\newcommand{\nPReasons}{.117}
\newcommand{\nPRiyadh}{.401}
\newcommand{\nPYear}{<.001}
\newcommand{\nPaired}{13{,}924}
\newcommand{\nPairsAffirmed}{417}
\newcommand{\nPairsDisturbed}{135}
\newcommand{\nPairsTested}{876}
\newcommand{\nReversed}{736}
\newcommand{\nReversedShare}{4.78}
\newcommand{\nSubstituted}{478}
\newcommand{\nSubstitutedShare}{3.11}
\newcommand{\nUnclearOutcome}{1{,}449}
\newcommand{\nUnclearShare}{9.42}
\newcommand{\nVaried}{3}
\newcommand{\nWindowFrom}{1439}
\newcommand{\nWindowTo}{1444}
\newcommand{\nWroteAffirming}{26.2}
\newcommand{\nWroteNotAdmitting}{17.4}
\newcommand{\nWroteReversing}{34.3}
\newcommand{\nWroteSubstituting}{33.2}


\newif\ifanon
\anonfalse

\title{Affirmed on the Reasons Below:\\
Appellate Review in \nAppeals{} Published Saudi Commercial Judgments}

\ifanon
  \author{}
  \date{}
\else
  \author{Abdullah Almohammedi \\
    Independent Researcher \\
    \texttt{abdullah.m.almohammedi@gmail.com} \\
    ORCID: \href{https://orcid.org/0009-0001-0832-0995}{0009-0001-0832-0995}}
  \date{}
\fi

\begin{document}
\doublespacing
\maketitle

\begin{abstract}
\noindent
Saudi commercial judgments are published with the appellate decision attached
to the judgment it reviews, giving \nPaired{} matched pairs of the same
dispute at two levels. \nNoReasonsShare{} per cent of appellate judgments give
no reasons of their own, affirming on the reasons below; whether a circuit
writes is not independent of what it decides, rising from
\nWroteAffirming{} per cent when it affirms to \nWroteReversing{} when it
reverses. Of appeals reaching the merits, \nDisturbedShare{} per cent disturb
the judgment. A logistic model on the \nModelN{} decided between
\nWindowFrom{} and \nWindowTo{} --- the window in which the record is mature
enough to read --- finds the procedural
share of a judgment's own citations the strongest protection against
disturbance (odds ratio \nOrProcedural{}), while its length, its citation
count and whether it wrote reasons predict nothing.
\end{abstract}

\noindent\textbf{Keywords:} appellate review; judicial reason-giving; empirical
legal studies; court administration; Saudi Arabia

\section{Introduction}

An appeal is studied where appellate decisions are digitised, which in
practice has meant the United States and a handful of other common-law
systems. The second instance elsewhere is known largely through court
statistics: how many appeals were brought, how many succeeded. What the
appellate court actually wrote, set beside what the court below wrote in the
same dispute, is rarely available at scale in any jurisdiction, because the
two documents live in different places.

In Saudi Arabia they live in the same record. The Ministry of Justice
publishes commercial judgments in full text, and where a case went up, the
appellate decision is published attached to the judgment it reviews. That
yields \nPaired{} matched pairs --- the same dispute seen twice --- inside a
corpus of \nJudgments{} judgments, with a further \nAppealOnly{} appellate
decisions published alone. Nothing about the pairing was designed for
research; it is a by-product of how the publisher stores a case file. But it
removes the objection that defeats most comparisons between levels, which is
that the cases differ. Here they do not differ. They are the same cases.

Three things follow, and the first is about the institution rather than any
case. \nNoReasonsShare{} per cent of appellate judgments in this corpus give
no reasons of their own: they affirm on the reasons below, in the formula
\emph{mahmulan `ala asbabihi}. The second instance therefore adds its own
statement of reasoning to the public record in about one appeal in four.
Whether it does so is not independent of what it decides --- circuits wrote
their own reasons in \nWroteAffirming{} per cent of affirmances and
\nWroteReversing{} per cent of reversals --- which matters both as a finding
and as a constraint on the third result below.

Second, of the appeals that reached the merits, \nDisturbedShare{} per cent
disturbed the judgment: reversed it, replaced it with a fresh judgment, or
varied it. The rest affirmed. A further \nNotAdmittedShare{} per cent of all
appellate decisions never reached the merits at all, failing at the threshold,
and are excluded from that rate rather than counted as affirmances, because an
inadmissible objection says nothing about whether the judgment below was
sound.

Third, a logistic model asks what separates the judgments that were disturbed
from those that survived, using only what can be read from the first-instance
document. On the \nModelN{} decided between \nWindowFrom{} and \nWindowTo{},
the strongest protection is the procedural share of a judgment's own
statutory citations (odds ratio \nOrProcedural{}, $p$ \nPProcedural{}). A
judgment resting on the law of procedure is harder to disturb than one
resting on the law of obligations. Its length, the number of instruments it
cites, and whether it wrote its own reasons predict nothing at all.

Two methodological points are reported rather than assumed. The appellate
disposition is read from the operative part of the appellate document and
nowhere else, because an appellate recital narrates the history of the case
and a pattern applied to the whole document reads an earlier affirmance as
this court's ruling; the classifier was corrected four times against
hand-read samples before it was used. And the year in which a judgment was
given cannot be read across the whole span: a recent judgment appears here
with an appellate decision only if the appeal was both brought and decided
before the corpus was collected, so the disturbed rate falls from
\nDisturbedShare{} per cent in the mature years to almost nothing in the last
of them. The window is therefore an argument to the model rather than a
silent assumption, and the article reports the mature one.

\section{Related work}
\label{sec:related}

\paragraph{What an appeal is for.} The economic account treats the appeal as
a device for correcting error at a cost, in which the party who lost has
private information about whether the decision was wrong and the system
relies on that information to select the cases worth reviewing (Shavell
1995). The account predicts a low disturbance rate among the appeals that are
brought, since litigants who know they were rightly beaten should not appeal;
what it does not predict is which surviving judgments differ from those
overturned, which is an empirical question this corpus can answer directly.

\paragraph{Measuring appellate outcomes.} The empirical literature on appeal
rates and outcomes is built largely on United States federal and state court
statistics, and is centrally concerned with who wins on appeal and why
plaintiffs do worse there than at trial (Eisenberg 2004). That work reads
outcomes from docket records: the appellate decision is a category rather
than a text. This article reads the text, which permits questions a docket
cannot answer --- whether the appellate court reasoned at all, and what it
cited when it did --- at the price of being confined to what one publisher
publishes.

\paragraph{Reason-giving, and its limits.} That courts should give reasons is
among the least controversial propositions in procedural theory, and among
the least absolute in practice. Cohen (2015) sets out the values that pull
against it --- efficiency, sincerity, and the risk of guidance a court is not
ready to give --- and shows common-law and civil-law systems converging on
similar compromises. Affirming on the reasons below is one such compromise,
and a defensible one: reasons adopted are reasons given. Whether a court
adopts or writes is nevertheless observable, and the rate at which the second
instance writes has not been measured for this jurisdiction. It is measured
here.

\paragraph{The filters between a dispute and this dataset.} Two selection
problems bound everything below and cannot be removed by better extraction.
Cases that settle never produce a judgment, and settled disputes differ
systematically from tried ones (Priest and Klein 1984); of the judgments
given, publication is a second filter, and published and unpublished
decisions have been shown to differ in kind rather than merely in number
(Siegelman and Donohue 1990). A third filter is specific to appeals and is
the more binding here: a judgment enters this analysis only if it was
appealed. Nothing in what follows describes judgments that no one challenged.

\paragraph{The jurisdiction.} Scholarship on Saudi law has been overwhelmingly
doctrinal and institutional, and until recently had little else to work with;
Vogel (2000) remains the standard English-language treatment of the courts,
written before judgments were published at scale. The commercial courts whose
work this article measures were reorganised under a new law in 2020, and the
corpus is dominated by the years since.

\section{Data}
\label{sec:data}

\subsection{A record that carries both levels}

\subsection{What the published record is}

\section{Method}
\label{sec:method}

\subsection{Reading the appellate disposition}

\subsection{Whether the appeal reasoned}

\subsection{Modelling disturbance}

\section{Results}
\label{sec:results}

\subsection{What appeals do}

\subsection{Reason-giving at the second instance}

\subsection{Reasoning compared, where both levels reasoned}

\subsection{What predicts disturbance}

\section{Threats to validity}
\label{sec:threats}

\section{Discussion}
\label{sec:discussion}

\section{Conclusion}

\section*{References}

\begin{hangparas}{1.5em}{1}
Cohen, Mathilde. 2015. ``When Judges Have Reasons Not to Give Reasons: A
Comparative Law Approach.'' \emph{Washington and Lee Law Review} 72(2): 483.

Eisenberg, Theodore. 2004. ``Appeal Rates and Outcomes in Tried and Nontried
Cases: Further Exploration of Anti-Plaintiff Appellate Outcomes.''
\emph{Journal of Empirical Legal Studies} 1(3): 659--688.

Priest, George L., and Benjamin Klein. 1984. ``The Selection of Disputes for
Litigation.'' \emph{Journal of Legal Studies} 13(1): 1--55.

Shavell, Steven. 1995. ``The Appeals Process as a Means of Error
Correction.'' \emph{Journal of Legal Studies} 24(2): 379--426.

Siegelman, Peter, and John J. Donohue III. 1990. ``Studying the Iceberg from
Its Tip: A Comparison of Published and Unpublished Employment Discrimination
Cases.'' \emph{Law \& Society Review} 24(5): 1133--1170.

Vogel, Frank E. 2000. \emph{Islamic Law and Legal System: Studies of Saudi
Arabia}. Leiden: Brill.
\end{hangparas}

\section*{Availability of data and code}

\ifanon\else
\section*{Declarations}
\fi

\end{document}
